% !TEX root = ../spec_kernel_agent_main.tex
\section{Introduction}
\label{sec:intro}

GPU kernels underpin every layer of modern LLM training and inference.
Even a few percent improvement to attention, GEMM, or normalization kernels translates directly into reduced infrastructure cost~\cite{kevin32b,cudaagent,cudaforge}.
However, kernel optimization is time and labor intensive, costing experts weeks of tuning kernel performance on specific GPUs~\cite{flashattention1,flashattention2,flashattention3}.
Agentic kernel optimization automates the manual kernel tuning procedure with reasoning LLMs. 
It casts the kernel optimization process as a feedback-guided iterative search over a large kernel design space~\cite{K-Search,Kernel-Smith,AlphaEvolve,cudaforge}.
Each iteration performs three phases in sequence:
\textbf{(1)} \emph{generation}: calling reasoning LLMs to generate candidate kernels;
\textbf{(2)} \emph{validation}: compiling and running these kernels against a reference kernel to check correctness and a first-cut speedup;
and \textbf{(3)} \emph{profiling}: measuring detailed performance metrics and adding them to the LLM's context to inform subsequent generations.
The objective of agentic kernel optimization is to \emph{maximize kernel performance under a fixed time or token budget}.

However, the large kernel design space makes this search time- and resource-costly, with
a kernel optimization task often taking tens of hours or even days~\cite{cudaforge,AlphaEvolve,kernelagent2025}.
To understand where the budget goes, we conduct a comprehensive workload characterization (see \S\ref{sec:mot}) of agentic kernel optimization across
10 \texttt{KernelBench}~\cite{kernelbench} tasks on an H200 testbed with two state-of-the-art open-source reasoning LLMs, 
GLM-5.1~\cite{glm5} and DeepSeek-V4-Pro~\cite{deepseekv4}.
Our workload characterization identifies three system-level inefficiencies that limit search efficiency: 
\textbf{(C1)} long generation latency: the generation phase accounts for 70--99\% of iteration time, directly capping the iteration count under fixed time budget;
\textbf{(C2)} insufficient profiling feedback: about 59--64\% of iterations consume GPU hours without producing profiling feedback to drive the next iteration's optimization; 
and \textbf{(C3)} underutilized validation and profiling resources: validation
and profiling GPUs achieve only 4.2--17.6\% utilization.
They remain idle during the generation phase because candidate kernels emerge
only after generation completes, and many candidates fail validation before profiling.

To address these inefficiencies, prior systems adopt different harness engineering
strategies, such as best-of-$K$~\cite{cudaforge,K-Search}, multi-agent~\cite{kernelagent2025}, 
and evolutionary algorithms~\cite{AlphaEvolve,Kernel-Smith}.
Despite their strategies, these systems share a common pattern. 
They dispatch multiple reasoning generations per iteration to increase the
candidate count and produce more profiling feedback, but this roughly multiplies
token cost by $K$. 
A strawman alternative is to dispatch non-reasoning generations to reduce token cost. 
However, non-reasoning generations rarely produce valid or performant kernels (see Table~\ref{tab:nonreasoning_speedups_w_wo_shared_prefix_cache}).
Moreover, prior systems 
do not reduce generation latency since candidate kernels become actionable only after the 
generation finishes, and validation/profiling resources remain idle during the reasoning 
process of the generation phase.

Our key insight is that the reasoning process exposes a window for producing
additional candidate kernels before the reasoning generation completes
(\S\ref{subsec:design_insight}).
Since non-reasoning generations are faster than reasoning generations, the system can dispatch
multiple non-reasoning generations in parallel with the ongoing reasoning generation to 
increase the candidate kernel count.
To improve the kernel performance of non-reasoning generations, we condition non-reasoning
generations on the prefix of the reasoning output.
With this conditioning, the non-reasoning generations can produce valid and performant kernels before the
reasoning generation completes (see Table~\ref{tab:nonreasoning_speedups_w_wo_shared_prefix_cache}).
Our further experiments show that prefix conditioning can even generate kernels with higher speedups than 
the average speedup of previously generated kernels (see Figure~\ref{fig:spec_reach95_token_cdf}).
This provides three benefits:
\textbf{(1)} shorter generation latency by terminating the reasoning
process early when a satisfactory candidate emerges (\textbf{C1});
\textbf{(2)} more profiling feedback through generating more candidate kernels per
iteration (\textbf{C2});
\textbf{(3)} higher validation/profiling utilization by keeping validation/profiling
resources busy during generation (\textbf{C3}).
We empirically show that a pragmatic termination criterion reduces generation
latency without compromising kernel performance
(see \S\ref{subsec:eval_termination_analysis}).
Building on this insight, we introduce \emph{speculative generation}, which
forks non-reasoning generations conditioned on the reasoning prefix.
Because each fork shares the cached reasoning prefix, speculative generation
incurs low token overhead (see \S\ref{subsec:eval_token_consumption}).
Moreover, speculative generation operates at the iteration level and is
orthogonal to token-level speculative decoding~\cite{specdec,medusa}.
The two techniques can be combined to further improve system efficiency.

However, realizing speculative generation introduces two system-level challenges.
\textbf{First,} \emph{how} to fork and \emph{when} to terminate:
the system must identify useful conditioning context from a noisy reasoning
trace and decide how many speculative generations to launch.
Too few forks leave validation/profiling resources idle, while too many overload
the queues and delay profiling feedback.
Terminating too early sacrifices kernel performance, while terminating too late saves negligible generation time.
\textbf{Second,} \emph{how} to manage bursty speculative-generation load.
The legacy static ``one GPU per kernel'' partitioning cannot handle this burstiness.
It leaves some GPUs idle while others queue, so resource allocation must adapt as bursts arrive.
The system also needs request prioritization to bound profiling feedback latency
and memory management to support bursty speculative generations in local LLM
deployments.

We present \sysname{}, an agentic kernel-optimization system that wraps a user-specified LLM, prompt, and search algorithm with two cooperating components.
\textbf{(1) SpecController} monitors the reasoning model's output stream and parses it to detect trigger signals, such as kernel-design decisions, fenced code blocks, and closed kernel bodies.
On a clean trigger, SpecController concatenates the iteration's prompt with the reasoning prefix and dispatches $K$ non-reasoning speculative generation requests.
The number of forks adapts to the available validation/profiling resource capacity, exposing enough candidates to keep GPUs busy without overloading the queues.
Once a speculatively generated kernel exceeds the average speedup of previously generated kernels, \sysname{} terminates the reasoning generation early to reduce the generation latency.
This threshold avoids terminating on weak speculative kernels while still allowing frequent early terminations.
\textbf{(2) ElasticScheduler} manages one elastic GPU pool, dynamically split between validation and profiling.
Each request takes any free GPU in its pool, so no GPU idles while another queues up.
Between iterations, ElasticScheduler reallocates the two pools based on the arrival rates of the previous iteration.
Within an iteration, the validation queue is served last-arrival-first.
Later candidates carry more reasoning prefix and are more likely to validate.
The profiling queue is served FIFO, so any validated kernel returns feedback fast and SpecController decides on the freshest signal.
We also observe that validation/profiling GPUs retain substantial unused memory
even under bursty speculative load.
We therefore repurpose this spare memory as remote storage for the
reasoning prefix cache, eliminating prefix recomputation of speculative generations.

Across \texttt{KernelBench} Level~1/2/3 tasks with two reasoning LLMs on H200,
\sysname{} consistently improves the efficiency of agentic kernel optimization
against CudaForge~\cite{cudaforge}, AlphaEvolve~\cite{AlphaEvolve}, and KernelAgent~\cite{kernelagent2025}.
Compared with these baselines, \sysname{} reduces end-to-end execution time by
1.68--1.82$\times$ and increases profiling feedback by 1.58--1.98$\times$.
It also improves final kernel speedup by 1.24--1.91$\times$, showing that
shorter E2E time does not come at the cost of kernel performance.
Finally, \sysname{} increases validation/profiling utilization from
4.2--17.6\% to 88.2--96.1\%.

This paper makes the following contributions:
\begin{itemize}
    \item \textbf{Characterization.} We conduct a comprehensive workload characterization of agentic kernel optimization and identify three system-level inefficiencies.
    \item \textbf{Insight.} We empirically establish that the reasoning process exposes a window for
    speculative generations conditioned on the reasoning prefix.
    \item \textbf{System.} We implement \sysname{}, comprising a SpecController and an ElasticScheduler that realize the above insight to address the three inefficiencies without sacrificing kernel performance.
    \item \textbf{Evaluation.} We evaluate \sysname{} on 20 \texttt{KernelBench} Level~1/2/3 tasks across two reasoning LLMs on H200 against three state-of-the-art agentic kernel optimization systems, with detailed results in \S\ref{sec:eval}.
\end{itemize}
